% GENERATED FILE. Edit canonical lessons, then run npm run generate:books.
% canonical-source-hash: fnv1a64:2288f99f73925cfd
% canonical-lessons: GE-C254-ich-finde-dass, GE-C254-meiner-meinung-nach, GE-C254-darum, GE-C254-review

\chapter{What I Think — ich finde, dass; darum}
\label{ch:ge-what-i-think}
\hlchaptermodality{254}

\begin{chapteropening}
\textbf{By the end of this chapter:} \emph{I can give an opinion with ich finde, dass or meiner Meinung nach, and draw a consequence with darum.}

Complete GE-C254-review: an opinion with dass and the verb at the end, a reason with weil, and a consequence with darum.
\end{chapteropening}

\section[ich finde, dass …]{ich finde, dass — an opinion, with its verb at the end}

\label{lesson:GE-C254-ich-finde-dass}



\begin{quote}
Say \textquotedblleft{}because German is beautiful\textquotedblright{}. (\emph{Weil Deutsch schön ist.}) Where did \emph{ist} go? (To the end.)
\end{quote}

\begin{grammarlens}[title={Grammar Lens: dass does what weil does}]
\textbf{dass} means \textquotedblleft{}that\textquotedblright{}, the \textquotedblleft{}that\textquotedblright{} of \textquotedblleft{}I think that …\textquotedblright{}. It opens a clause, and like \emph{weil} it sends the verb to the end:

\begin{quote}
\emph{Ich finde, dass Deutsch schön} \textbf{ist}. — I think German is beautiful. \emph{Ich glaube, dass der Zug Verspätung} \textbf{hat}. — I believe the train is late.
\end{quote}

German sets the \emph{dass} clause off with a comma. German can drop \emph{dass}, as English drops \textquotedblleft{}that\textquotedblright{}, but then the verb returns to second place: \emph{Ich glaube, der Zug} \textbf{hat} \emph{Verspätung.} Both are standard.

To disagree, put \emph{nicht} on the opinion verb: \emph{Ich finde nicht, dass das richtig ist.}
\end{grammarlens}

\begin{culture}
\textbf{ich finde} is a judgement of your own, often about something you have tried.

\textbf{ich glaube} is a belief you hold without being sure.

\textbf{ich meine} is \textquotedblleft{}I think, I mean\textquotedblright{}. The verb is \textbf{meinen}, and you know its noun already: \emph{die Meinung}, the opinion. \emph{Meinen} also means \textquotedblleft{}to mean\textquotedblright{}: \emph{Was meinst du?}, what do you mean?

One trap in the spelling: \textbf{dass} with two \emph{s} joins clauses; \textbf{das} with one is \textquotedblleft{}the\textquotedblright{} or \textquotedblleft{}that one\textquotedblright{}. Before 1996 the joiner was written \emph{daß}. Swiss spelling does without \emph{ß} altogether, so Swiss writers used \emph{dass} long before then.
\end{culture}

\begin{cousinweb}
\textbf{dass} and \textbf{das} go back to one Old High German word, \emph{daz}, from Proto-Germanic \emph{*þat} (EtymDB 2.1), which EtymDB sets beside English \emph{that}. English kept one spelling for both jobs; German later split it in two, by the usual account in the sixteenth century.

\textbf{meinen} $\leftarrow$ Old High German \emph{meinen} $\leftarrow$ Proto-Germanic \emph{*mainijaną}, \textquotedblleft{}to mean, to think\textquotedblright{} (EtymDB 2.1): the same verb as English \emph{mean}.
\end{cousinweb}

\subsection*{Your turn}
Say these aloud:

\begin{itemize}
\raggedright
  \item \emph{Ich finde, dass Deutsch schön ist.}
  \item \emph{Ich glaube, dass der Zug Verspätung hat.}
  \item the same without \emph{dass}: \emph{Ich glaube, der Zug hat Verspätung.}
  \item \emph{Ich meine, dass das richtig ist.}
  \item \emph{Ich finde nicht, dass das richtig ist.}
\end{itemize}

\begin{tcolorbox}[breakable,colback=teal!4,colframe=teal!35!black,title={Before you move on}]
After \emph{dass}, where does the verb go? (\textbf{To the end.}) Drop \emph{dass}, and where does it go? (\textbf{Back to second place.}) Which opinion verb is a judgement you formed yourself? (\emph{\textbf{Finden.}}) And which noun belongs to \emph{meinen}? (\emph{\textbf{Die Meinung.}})
\end{tcolorbox}
\section[meiner Meinung nach]{meiner Meinung nach — in my opinion}

\label{lesson:GE-C254-meiner-meinung-nach}



\begin{quote}
Say \textquotedblleft{}I think that is right\textquotedblright{}, with \emph{dass}. (\emph{Ich finde, dass das richtig ist.})

There is a shorter way in, and it does not move the verb to the end.
\end{quote}

\subsection*{You'll want to know: meiner Meinung nach}
\textbf{meiner Meinung nach} — \textquotedblleft{}in my opinion\textquotedblright{}: \emph{MY-ner MY-noong nahkh}.

\begin{quote}
\emph{Meiner Meinung nach ist das richtig.} — In my opinion, that is right.
\end{quote}

Two things are unusual in it. \textbf{Nach} comes after its noun, and here it means \textquotedblleft{}according to\textquotedblright{}. \emph{Nach meiner Meinung} exists too, but \emph{nach} last is the usual order. And \textbf{meiner} is \emph{mein} with the dative ending \emph{-er}, the case \emph{nach} asks for. With \textquotedblleft{}your\textquotedblright{}: \emph{deiner Meinung nach}.

\begin{grammarlens}[title={Grammar Lens: one place, then the verb}]
The whole phrase, three words long, takes the first place of the sentence, exactly as \emph{dann} did. So the verb comes next, and the subject waits behind it:

\begin{quote}
\emph{Meiner Meinung nach} \textbf{ist} \emph{das wichtig.} \emph{Meiner Meinung nach} \textbf{hat} \emph{der Zug Verspätung.}
\end{quote}

Compare the two ways in:

\begin{quote}
\emph{Ich finde, dass das wichtig} \textbf{ist}. — the verb at the end. \emph{Meiner Meinung nach} \textbf{ist} \emph{das wichtig.} — the verb second.
\end{quote}

Same opinion, two word orders. \emph{Meiner Meinung nach} is at home in a discussion and in writing.
\end{grammarlens}

\begin{cousinweb}
\textbf{Meinung} is the verb \emph{meinen} + \emph{-ung}, from Old High German \emph{meinen} (EtymDB 2.1).

\textbf{nach} is set by EtymDB 2.1 beside \emph{nah}, \textquotedblleft{}near\textquotedblright{}, from Old High German \emph{nāh}, and beside English \emph{nigh}. From \textquotedblleft{}near\textquotedblright{} it came to mean \textquotedblleft{}after\textquotedblright{} and, behind a noun, \textquotedblleft{}according to\textquotedblright{}; EtymDB does not record those steps, so they are given here as the usual account.
\end{cousinweb}

\subsection*{Your turn}
Say these aloud:

\begin{itemize}
\raggedright
  \item \emph{meiner Meinung nach}
  \item \emph{Meiner Meinung nach ist das richtig.}
  \item \emph{Meiner Meinung nach ist das wichtig.}
  \item the same opinion with \emph{dass}: \emph{Ich finde, dass das wichtig ist.}
  \item to a friend: \emph{Deiner Meinung nach?}
\end{itemize}

\begin{tcolorbox}[breakable,colback=teal!4,colframe=teal!35!black,title={Before you move on}]
\textquotedblleft{}In my opinion, that is important\textquotedblright{}? (\emph{\textbf{Meiner Meinung nach ist das wichtig.}}) Where does \emph{nach} stand? (\textbf{After its noun.}) After the phrase, which comes next, the verb or the subject? (\textbf{The verb.})
\end{tcolorbox}
\section[darum]{darum — that's why}

\label{lesson:GE-C254-darum}



\begin{quote}
Say the German for \textquotedblleft{}therefore\textquotedblright{}, from chapter two hundred and fifty. (\emph{Deshalb.}) And the question \textquotedblleft{}why?\textquotedblright{} (\emph{Warum?})
\end{quote}

\subsection*{You'll want to know: darum}
\textbf{darum} — \textquotedblleft{}that's why, so\textquotedblright{}: \emph{da-RUM}, with a short \emph{u}; when you point hard at the reason, \emph{DA-rum}.

\begin{quote}
\emph{Es regnet, darum nehme ich den Regenschirm.} — It's raining, so I'm taking the umbrella. \emph{Der Zug hatte Verspätung, darum kam ich spät.} — The train was late, that's why I came late.
\end{quote}

\emph{Darum} and \emph{deshalb} do the same job, and either fits in both sentences. \emph{Darum} has one more use: on its own, it is the curt answer to \emph{warum}. \textquotedblleft{}\emph{Warum?}\textquotedblright{} \textquotedblleft{}\emph{Darum!}\textquotedblright{}, because.

\begin{grammarlens}[title={Grammar Lens: two directions, two word orders}]
One fact, told two ways:

\begin{quote}
\emph{Ich nehme den Regenschirm, weil es} \textbf{regnet}. \emph{Es regnet, darum} \textbf{nehme} \emph{ich den Regenschirm.}
\end{quote}

\textbf{weil} points back to the cause, and its verb goes to the end. \textbf{darum} points forward to the result. It takes the first place of its clause, like \emph{dann}, so its verb comes straight after it: \emph{darum nehme ich}, not \emph{darum ich nehme}.

The comma before \emph{darum} is where one clause stops and the next begins. In a longer stretch, \emph{darum} can open a fresh sentence: \emph{Es regnet. Darum nehme ich den Regenschirm.}
\end{grammarlens}

\begin{cousinweb}
\textbf{darum} is \emph{da(r)}, \textquotedblleft{}there, that\textquotedblright{}, + \textbf{um}, \textquotedblleft{}around, about\textquotedblright{}: about that. You took \emph{warum} apart the same way, \textquotedblleft{}about what\textquotedblright{}. EtymDB 2.1 lists \emph{darum} without its parts, but traces \emph{um} through Old High German \emph{umbi} to Proto-Germanic \emph{*umbi}, beside Greek \emph{amphí} and Latin \emph{ambi-}, \textquotedblleft{}around\textquotedblright{}.
\end{cousinweb}

\subsection*{Your turn}
Say these aloud:

\begin{itemize}
\raggedright
  \item \emph{darum}
  \item \emph{Es regnet, darum nehme ich den Regenschirm.}
  \item the same fact backwards: \emph{Ich nehme den Regenschirm, weil es regnet.}
  \item \emph{Der Zug hatte Verspätung, darum kam ich spät.}
  \item the same with \emph{deshalb}: \emph{Der Zug hatte Verspätung, deshalb kam ich spät.}
\end{itemize}

\begin{tcolorbox}[breakable,colback=teal!4,colframe=teal!35!black,title={Before you move on}]
\textquotedblleft{}So I'm taking the umbrella\textquotedblright{}? (\emph{\textbf{Darum nehme ich den Regenschirm.}}) Which word points back to the cause? (\emph{\textbf{Weil.}}) Which points forward to the result? (\emph{\textbf{Darum}}, or \emph{deshalb}.) And which little word is in both \emph{warum} and \emph{darum}? (\emph{\textbf{Um}}, around.)
\end{tcolorbox}
\section[(review)]{Review — what I think, and why}

\label{lesson:GE-C254-review}



\begin{quote}
Say \textquotedblleft{}I believe that the train is late\textquotedblright{}. (\emph{Ich glaube, dass der Zug Verspätung hat.})
\end{quote}

\begin{grammarlens}[title={What you've built}]
Three lessons, and an argument has all its joints. An opinion: \emph{ich finde}, \emph{ich glaube} or \emph{ich meine}, then \textbf{dass} and the verb at the end. Or \textbf{meiner Meinung nach} in the first place, with the verb right after it. A reason with \textbf{weil}, verb last. A consequence with \textbf{darum}, verb second.

Two sentences, and both word orders in them:

\begin{quote}
\emph{Meiner Meinung nach} \textbf{kommen} \emph{wir spät, weil der Zug Verspätung} \textbf{hat}. \emph{Der Zug hat Verspätung, darum} \textbf{nehme} \emph{ich ein Taxi.}
\end{quote}

Verb second after the opinion phrase, verb last after \emph{weil}, and verb second again after \emph{darum}.
\end{grammarlens}

\subsection*{Your turn}
Say these aloud:

\begin{itemize}
\raggedright
  \item \emph{Ich finde, dass Deutsch schön ist.}
  \item \emph{Meiner Meinung nach ist das richtig.}
  \item \emph{Ich nehme den Regenschirm, weil es regnet.}
  \item \emph{Es regnet, darum nehme ich den Regenschirm.}
  \item one opinion of your own with \emph{ich finde, dass}, and one consequence with \emph{darum}
\end{itemize}

\begin{tcolorbox}[breakable,colback=teal!4,colframe=teal!35!black,title={Before you move on}]
After \emph{dass} and \emph{weil}, where is the verb? (\textbf{At the end.}) After \emph{darum} and \emph{meiner Meinung nach}? (\textbf{Second}, straight after them.) Which one points back to a cause? (\emph{\textbf{Weil.}})
\end{tcolorbox}
